
Large language model deployment for applications requiring reliability faces a trade-off between computational efficiency and uncertainty quantification quality. Zero-cost uncertainty methods offer efficiency, while expensive methods such as Monte Carlo dropout impose 5-$10\times$ inference overhead. Current UQ research reports performance rankings without cost analysis, providing limited guidance for practitioners operating under budget constraints. A system with $5\times$ cost tolerance cannot determine whether MC dropout k=5 justifies overhead versus temperature scaling at $0\times$ cost.

Existing work evaluates UQ methods in isolation or compares methods without measuring computational cost. No study maps the complete cost-performance space on the same selective prediction task. Research focuses on identifying the best-performing method rather than characterizing trade-offs, treating cost as secondary rather than as a deployment constraint. Practitioners consequently lack information for matching UQ method selection to deployment budgets.

The gap is concrete: no systematic benchmark compares AUROC versus inference cost across UQ method variants on selective prediction. Cost reporting varies across studies, k-values differ, FLOPs are rarely measured, and winner-take-all evaluation obscures trade-off information.

Our observation is that cost-performance trade-offs exist because different UQ mechanisms operate in distinct efficiency zones. If multiple methods occupy the Pareto frontier where no method achieves both higher AUROC and equal-or-lower cost, practitioners can select methods based on budget. Temperature scaling at $1\times$ cost versus MC dropout k=5 at $5\times$ cost: neither dominates if temperature scaling achieves lower AUROC at zero overhead.

Our contributions are:

\textbf{Systematic cost-performance benchmark for UQ on LLM selective prediction.} We compare 6 UQ method variants (temperature scaling, conformal prediction, MC dropout k=1/3/5/10) on TruthfulQA selective prediction (817 questions), measuring AUROC and FLOPs-normalized inference cost. This Pareto frontier framing reveals the trade-off space rather than declaring a single winner.

\textbf{Efficiency point for MC dropout.} MC dropout k=3 achieves AUROC 0.704 (exceeding 0.70 threshold) at $3\times$ cost. This contrasts with prior work using k=30 or higher, indicating that Bayesian approximation converges at lower k for 8B model scale.

\textbf{Epistemic uncertainty threshold at 8B scale.} Zero-cost methods (temperature scaling 0.682, conformal prediction 0.695) fall below the 0.70 AUROC threshold, while MC dropout $k\geq 3$ exceeds it. This indicates that post-hoc calibration is insufficient for selective prediction at 8B scale and that epistemic uncertainty via stochastic forward passes is required.

\textbf{Pareto-optimal methods across cost zones.} Results show temperature scaling, conformal prediction, and MC dropout k=3/5/10 all occupy the Pareto frontier, confirming no universal dominance. This validates budget-aware UQ selection as a framework.

Section 2 reviews UQ methods and identifies the missing cost-performance analysis. Section 3 presents our Pareto frontier methodology. Section 4 describes experimental setup on TruthfulQA. Section 5 presents results showing Pareto-optimal methods. Section 6 discusses epistemic versus aleatoric uncertainty at 8B scale. Section 7 concludes.
